\chapter*{Declaration on the Use of Generative AI}
\addcontentsline{toc}{chapter}{Declaration on the Use of Generative AI}

In accordance with the guidelines for the responsible use of AI, the author acknowledges the use of Generative AI tools
(specifically Google Gemini) to assist with English language proofreading, sentence structure optimization, and the
drafting and debugging of the embedded firmware code.

The author certifies that the conceptual design and scientific core of the thesis, including the mechanical design of
the hexapod limb, the custom PCB design, and the interpretation of experimental results, are the result of the author's
original work.

All content and code segments generated or suggested by AI tools were subjected to a rigorous critical review and
verification by the author to ensure accuracy and compliance with the project specifications. The author bears full
responsibility for the final content of this thesis.


\chapter*{Ringraziamenti}
\addcontentsline{toc}{chapter}{Ringraziamenti}

Chi mi conosce sa che solitamente preferisco esprimere affetto o gratitudine attraverso gesti piuttosto che a parole.
Molto spesso la causa è la vergogna o l'imbarazzo, che mi fanno sentire vulnerabile, ma per questa occasione ho deciso
di fare uno sforzo.

Un grazie speciale ai miei genitori, Enzo e Manuela, da sempre i miei sostenitori numero 1. Mi sopportano dal mio primo
giorno di vita (un lavoro molto duro, ne sono consapevole). Sappiate che anche se non sono bravo a farvelo capire, vi
voglio un mondo di bene e vi porterò sempre nel mio cuore.

Ringrazio mia sorella Annalaura, che seppur controvoglia, alla fine si impegna sempre per darmi una mano quando serve.

Ringrazio Letizia, la mia compagna, con cui ormai da anni condivido passioni ed avventure. Sono sicuro che il futuro ci
riserverà molti altri momenti speciali.

Ringrazio Mattia, Linda e Marco M., i miei migliori amici. Con voi riesco sempre ad essere felice e abbandonare ogni
preoccupazione, anche nei momenti più bui.

Grazie a tutti gli amici, che negli anni mi hanno accompagnato, tra alti e bassi. In particolare il gruppo del Fantabosco,
(Max, Matteo B., Lucrezia, Fra e Daniel) con cui ho condiviso serate indimenticabili, gli ultimi arrivati dell'università
(Gio, Giulia, Liuk, Fede, Paolo, Anna, Edo, Noemi, Mary, RKB ed Elena) con cui sono sopravvissuto tra un esame e l'altro,
a questi ultimi tre anni, ma anche il gruppo della laurea triennale (Marco L., Matteo V., Tomas, Marta e Bias).
Tutti voi, sappiate che anche se dovessimo non rimanere più in contatto, da parte mia sarà come non avervi mai salutati,
e se avrete bisogno di aiuto saprete dove trovarmi.

Un sentito ringraziamento anche al Professor Angelo Cenedese, che mi ha permesso di proseguire questo progetto personale
fino a qui, mostrandosi da subito molto disponibile.

Infine, ringrazio chiunque abbia mai incrociato il suo cammino con il mio, nel bene e nel male. Sono convinto che senza
anche uno solo di voi, non sarei la persona che oggi sono fiero di essere.